\chapter{Kikuchi LDC — Two XOR}
%%% generated by the blueprint pipeline from TCSlib.KikuchiLDC.LDC.TwoXOR
\section{Overview}

This module covers \S7 of the Alrabiah--Guruswami--Kothari--Manohar argument: refuting
the $2$-XOR instance produced by the hypergraph decomposition. It records the abstract
value of the $2$-XOR polynomial attached to the bipartite matchings $G_i$, and the
refutation bound stating that this value is $O\!\left(n k \sqrt{\log n / d}\right)$ once
the number of heavy pairs satisfies $\abs{P} \le nk/d$.

\section{Declarations}

\begin{definition}[Value of the 2-XOR polynomial]
\lean{twoXORVal}
\label{twoXORVal}
\leanok
\statementsource{arXiv.2308.15403}{pdf-b372a20892ce-p006-b004}
For parameters $k, n, d$ and a message $b : \mathrm{Fin}\,k \to \mathbb{Z}$, the modelled
value of the $2$-XOR polynomial $g_b$ obtained from the decomposition, namely the real
number
\[
  \mathrm{val}(g_b) \;=\; n \, k \, \sqrt{\frac{\log n}{d}} .
\]
The definition is abstract: the quantity depends only on $n$, $k$ and the threshold $d$,
and is constant in the edge count and in $b$.
\end{definition}

\begin{theorem}[2-XOR refutation bound]
\lean{two_xor_refutation}
\label{two_xor_refutation}
\leanok
\difficulty{5}
\statementsource{arXiv.2308.15403}{
  pdf-b372a20892ce-p006-b004,
  pdf-b372a20892ce-p007-b002
}
Fix natural numbers $k, n, d$ with $k \ge 1$, $n \ge 2$, and $d \ge 1$. Then there is a
positive constant $C_0 > 0$ such that for every natural number $n_P$ satisfying $n_P\, d
\le n k$, one has
\[
  n \, \sqrt{n_P} \, \sqrt{k \log n}
    \;\le\; C_0 \, n \, k \, \sqrt{\frac{\log n}{d}} ,
\]
where $\log$ denotes the natural logarithm.
\end{theorem}
